% Part: incompleteness
% Chapter: representability-in-q
% Section: beta-function

\documentclass[../../../include/open-logic-section]{subfiles}

\begin{document}

\olfileid[hi]{inc}{req}{bet}
\olsection{बीटा फलन उपपत्ति}

यह दिखाने के लिए कि जोड़, गुणन और $\Char{=}$ उपलब्ध होने पर हम आदिम
पुनरावर्तन कर सकते हैं, हमें अनुक्रमों को सँभालने वाले फलन विकसित करने
होंगे। (यदि हमारे पास घातांक भी होता, तो हमारा कार्य अधिक सरल होता।) जब
हमारे पास आदिम पुनरावर्तन था, तब हम ``$n$-वीं अभाज्य संख्या'' जैसी चीजें
परिभाषित करके काफी सीधा कूटलेखन चुन सकते थे। किंतु यहाँ हमारे पास आदिम
पुनरावर्तन नहीं है---वास्तव में हम यह दिखाना चाहते हैं कि न्यूनीकरण का
उपयोग करके आदिम पुनरावर्तन किया जा सकता है---इसलिए हमें कुछ अधिक चतुर होना
होगा।

\begin{lem}
\ollabel{lem:beta}
ऐसा कोई फलन $\beta(d,i)$ है कि प्रत्येक अनुक्रम $a_0$, \dots,~$a_n$ के
लिए कोई संख्या~$d$ ऐसी है कि प्रत्येक $i \le n$ के लिए
$\beta(d,i) = a_i$। इसके अतिरिक्त $\beta$ को मूल फलनों से केवल संयोजन और
नियमित न्यूनीकरण द्वारा परिभाषित किया जा सकता है।
\end{lem}

$d$ को अनुक्रम $\tuple{a_0, \dots, a_n}$ का कूट और $\beta(d,i)$ को उसका
$i$-वाँ अवयव लौटाने वाला फलन समझें। (ध्यान दें कि इस ``कूटलेखन'' में वह
अभाज्य-घात कूटलेखन \emph{नहीं} प्रयुक्त होता जिससे हम पहले ही परिचित हैं!)
उपपत्ति बहुत न्यूनतम दावा करती है; वह यह नहीं कहती कि हम अनुक्रमों को जोड़
सकते हैं या उनमें अवयव अंत में लगा सकते हैं, बल्कि यह भी नहीं कहती कि
संयोजन और नियमित न्यूनीकरण से परिभाष्य फलनों द्वारा $a_0$, \dots,~$a_n$ से
$d$ को \emph{संगणित} कर सकते हैं। वह केवल इतना कहती है कि ऐसा एक
``विकूटन'' फलन है जिससे प्रत्येक अनुक्रम ``कूटबद्ध'' होता है।

$\beta$ संकेतन ग्योडेल (G\"odel) का है। फिर कहें तो उपपत्ति सिद्ध करने का
कठिन भाग इन देखने में सीमित साधनों से उपयुक्त~$\beta$ परिभाषित करना है,
अर्थात् केवल संयोजन और न्यूनीकरण से---यद्यपि हमें जोड़, गुणन और~$\Char{=}$
का उपयोग करने की अनुमति है। इस उपपत्ति को सिद्ध करने के कई ढंग हैं, किंतु
सबसे साफ ढंगों में अब भी ग्योडेल (G\"odel) की मूल विधि है, जिसमें सुन ज़ी प्रमेय
(परंपरागत नाम ``चीनी शेषफल प्रमेय'') नामक संख्या-सैद्धांतिक तथ्य का उपयोग
हुआ था।

\begin{defn}
दो प्राकृतिक संख्याएँ $a$ और $b$ \emph{परस्पर अभाज्य} तभी और केवल तभी हैं
जब उनका महत्तम समापवर्तक~$1$ हो; दूसरे शब्दों में, उनका कोई और उभयनिष्ठ
भाजक न हो।
\end{defn}

\begin{defn}
प्राकृतिक संख्याएँ $a$ और $b$, \emph{मापांक~$c$ के अनुसार सर्वांगसम} हैं,
$a \equiv b \mod c$, तभी और केवल तभी जब $c \mid (a-b)$; अर्थात्~$c$ से
भाग देने पर $a$ और $b$ का शेषफल समान हो।
\end{defn}

सुन ज़ी प्रमेय यह है:
\begin{thm}
मान लें कि $x_0$, \dots,~$x_n$ (युग्मशः) परस्पर अभाज्य हैं।
$y_0$, \dots,~$y_n$ कोई भी संख्याएँ हों। तब कोई संख्या $z$ ऐसी है कि
\begin{align*}
z & \equiv y_0 \mod x_0 \\
z & \equiv y_1 \mod x_1 \\
& \vdots  \\
z & \equiv y_n \mod x_n.
\end{align*}
\end{thm}

हम सुन ज़ी प्रमेय का उपयोग इस प्रकार करेंगे: यदि $x_0$, \dots,~$x_n$
क्रमशः $y_0$, \dots,~$y_n$ से बड़े हैं, तो हम $z$ को अनुक्रम
$\tuple{y_0, \dots,y_n}$ का कूट मान सकते हैं। $y_i$ को वापस पाने के लिए
केवल $z$ को~$x_i$ से भाग देकर शेषफल लेना होगा। इस कूटलेखन का उपयोग करने
के लिए हमें $x_0$, \dots,~$x_n$ के उपयुक्त मान खोजने होंगे।

इस संबंध में कुछ प्रेक्षण हमारी सहायता करेंगे। $y_0$, \dots,~$y_n$ दिए
जाने पर मानें
\begin{align*}
j &= \max(n, y_0 + 1, \dots, y_n + 1), \\
m &= \lcm(1,\dots,j),
\end{align*}
और मानें
\begin{align*}
x_0 & = 1 + m \\
x_1 & = 1 + 2 \cdot m \\
x_2 & = 1 + 3 \cdot m \\
& \vdots  \\
x_n & = 1 + (n+1) \cdot m.
\end{align*}
तब दो बातें सत्य हैं:
\begin{enumerate}
\item\ollabel{rel-prime} $x_0,\dots,x_n$ परस्पर अभाज्य हैं।
\item\ollabel{less} प्रत्येक $i$ के लिए $y_i < x_i$।
\end{enumerate}
यह देखने के लिए कि \olref{rel-prime} सत्य है, ध्यान दें कि यदि $p$ अभाज्य
संख्या है और $p \mid x_i$ तथा $p \mid x_k$, तो
$p \mid 1 + (i+1) m$ और $p \mid 1 + (k+1) m$। किंतु तब $p$ उनके अंतर को
विभाजित करता है,
\[
(1 + (i+1)m) - (1+ (k+1)m) = (i-k) m.
\]
चूँकि $p$,~$1 + (i+1)m$ को विभाजित करता है, वह~$m$ को भी विभाजित नहीं कर
सकता (अन्यथा पहले भाग में शेषफल~$1$ आता)। अतः $p$,~$i-k$ को विभाजित करता
है, क्योंकि $p$,~$(i-k)m$ को विभाजित करता है। किंतु
$\left|i-k\right|$ अधिकतम~$n$ है और हमने $j \geq n$ चुना है, इसलिए इससे
$p \mid m$ निकलता है, जो फिर विरोधाभास है। अतः कोई अभाज्य संख्या $x_i$
और $x_k$ दोनों को विभाजित नहीं करती। अनुच्छेद~\olref{less} सरल है: हमारे
पास $y_i < j \leq m < x_i$ है।

अब $\beta$ फलन उपपत्ति सिद्ध करें। याद करें कि हम $0$, उत्तराधिकारी, जोड़,
गुणन, $\Char{=}$, प्रक्षेप और उनसे संयोजन तथा नियमित फलनों पर लागू
न्यूनीकरण द्वारा परिभाषित किसी भी फलन का उपयोग कर सकते हैं। किसी संबंध का
अभिलाक्षणिक फलन इस प्रकार परिभाष्य हो तो हम उस संबंध का भी उपयोग कर सकते
हैं। पहले की तरह हम दिखा सकते हैं कि ये संबंध बूलीय संयोजनों और परिबद्ध
परिमाणीकरण के अंतर्गत संवृत हैं; उदाहरणतः:
\begin{align*}
\fn{not}(x) & \defis \Char{=}(x,0)\\
\bmin{x \leq z}{R(x,y)} & \defis \umin{x}{(R(x,y) \lor x = z)}\\
\bexists{x \leq z}{R(x,y)} & \defiff R(\bmin{x \leq z}{R(x,y)}, y).
\end{align*}
तब हम दिखा सकते हैं कि निम्न सभी भी आदिम पुनरावर्तन के बिना परिभाष्य हैं:
\begin{enumerate}
\item युग्मन फलन, $J(x,y) = \frac{1}{2}[(x+y)(x+y+1)] + x$;
% maybe explain more what is going on here, a bit confusing.
\item प्रक्षेप फलन
\begin{align*}
K(z) & = \bmin{x \leq z}{\bexists{y \leq z}{z = J(x,y)}},\\
L(z) & = \bmin{y \leq z}{\bexists{x \leq z}{z = J(x,y)}};
\end{align*}
\item लघुतर संबंध $x < y$;
\item भाज्यता संबंध $x \mid y$;
% \item $x \tsub y$
% \item $\fn{Prime}(x)$
% \item Assuming $p$ is prime, the relation ``$x$ is a power of $p$'':
% \[
% \bforall{y \leq x}{(y \mid x \lif y = 1 \lor y = x)}.
% \]
\item फलन $\fn{rem}(x,y)$, जो $y$ को~$x$ से भाग देने पर प्राप्त शेषफल
  लौटाता है।
\end{enumerate}
अब परिभाषित करें
\begin{align*}
\beta^*(d_0,d_1,i) & = \fn{rem}(1+(i+1) d_1,d_0) \text{ और}\\
\beta(d,i) & = \beta^*(K(d),L(d),i).
\end{align*}
यही वह फलन है जो हमें चाहिए। ऊपर की भाँति $a_0,\dots,a_n$ दिए जाने पर
मानें
\[
j = \max(n,a_0+1,\dots,a_n+1),
\]
और $d_1 = \lcm(1,\dots,j)$ मानें। ऊपर के \olref{rel-prime} के अनुसार हम
जानते हैं कि $1+d_1$, $1+2 d_1$, \dots, $1+(n+1) d_1$ परस्पर अभाज्य हैं,
और~\olref{less} के अनुसार वे क्रमशः $a_0,\dots,a_n$ से बड़े हैं। सुन ज़ी
प्रमेय के अनुसार कोई मान~$d_0$ ऐसा है कि प्रत्येक~$i$ के लिए
\[
d_0 \equiv a_i \mod (1+(i+1)d_1)
\]
और इसलिए ($d_1$,~$a_i$ से बड़ा होने के कारण)
\[
a_i = \fn{rem}(1+(i+1)d_1,d_0).
\]
$d = J(d_0,d_1)$ मानें। तब प्रत्येक $i \le n$ के लिए हमारे पास
\begin{align*}
\beta(d,i) & =  \beta^*(d_0,d_1,i) \\
& =  \fn{rem}(1+(i+1) d_1,d_0) \\
& =  a_i
\end{align*}
है, और यही हमें चाहिए था। इससे $\beta$-फलन उपपत्ति का प्रमाण पूरा होता है।

\begin{prob}
  दिखाएँ कि संबंध $x < y$, $x \mid y$ और फलन~$\fn{rem}(x,y)$ को आदिम
  पुनरावर्तन के बिना परिभाषित किया जा सकता है। आप $0$, उत्तराधिकारी, जोड़,
  गुणन, $\Char{=}$, प्रक्षेप तथा परिबद्ध न्यूनीकरण और परिमाणीकरण का उपयोग
  कर सकते हैं।
\end{prob}

\end{document}
